% Source: https://artofproblemsolving.com/community/c3h3129186_2023_midmichigan_mathematical_olympiad__grades_56
% Stable identity: aops-topic-3129186-q02
Reduce the pile of 4 to 2, leaving $(1,2,3)$. More generally, whenever at least two piles exceed 1, leave a position with bitwise XOR zero; whenever exactly one pile exceeds 1, reduce that pile to 0 or 1 so that the total number of piles of size 1 becomes odd. When all piles have size 1, an odd number is losing, since turns alternate and the player making the final move loses. To justify the XOR rule, a move changes the XOR to a nonzero value. From a nonzero XOR position, choose a pile carrying the highest nonzero bit of that XOR and replace its size $a$ by $a\mathbin{\mathrm{xor}}X<a$; this restores zero. A zero-XOR position cannot have exactly one pile larger than 1. If an opponent leaves exactly one large pile, apply the special rule above instead. Otherwise the zero-XOR restoration remains available. Thus the strategy preserves a losing position for the opponent until the final one-pile stage. Since $1\mathbin{\mathrm{xor}}2\mathbin{\mathrm{xor}}3=0$, the stated first move wins.
